\documentclass[../../../../main.tex]{subfiles}
\begin{document}

\begin{flushright}
\footnotesize\textit{【自動文字起こし・要確認】}
\end{flushright}

{\bf [解]}
$\triangle ABC$を$AB=AC=1$，$BC=a$（$0<a<2$），$\angle CAB=\alpha$，$\angle ABC=\beta$とする．周上の2点$P,Q$を結ぶ線分が$\triangle ABC$の面積を2等分するとき，$P,Q$の位置は対称性から次の2通りに帰着できる．

\textbf{$1^\circ$ $P,Q$が辺$AB,AC$上にある場合}

$AP=x,\ AQ=y$（$0<x,y\le1$）とおく．$\triangle APQ$が$\triangle ABC$の面積の半分の時，
\begin{align*}
xy=\frac12 \quad\cdots\text{①}
\end{align*}
$\triangle APQ,\ \triangle ABC$に余弦定理を用いて
\begin{align*}
PQ^2=x^2+y^2-2xy\cos\alpha,\qquad a^2=1+1-2\cos\alpha
\end{align*}
後者から$\cos\alpha=\dfrac{2-a^2}2$，これと①を前者に代入して
\begin{align*}
PQ^2=x^2+y^2-2xy\cdot\frac{2-a^2}2=x^2+y^2-\frac{2-a^2}2 \quad\cdots\text{②}
\end{align*}
（①より$2xy=1$を用いた．）①，$0<x,y\le1$のもとで$x^2+y^2$を最小化する．相加相乗平均から$x^2+y^2\ge2xy=1$，等号は$x=y$の時．①とあわせて$x=y=1/\sqrt2$（$\le1$をみたす）で最小値$x^2+y^2=1$をとる．②より
\begin{align*}
\min PQ^2=1-\frac{2-a^2}2=\frac{a^2}2
\end{align*}

\textbf{$2^\circ$ $P,Q$が辺$AB,BC$上にある場合}（$AC,BC$上の場合も対称性から同じ結論になる）

$BP=x$（$0<x\le1$），$BQ=y$（$0<y\le a$）とおく．$\triangle BPQ$が半分の面積の時，$\triangle ABC$の面積が$\frac12\cdot1\cdot a\sin\beta$であることから
\begin{align*}
xy=\frac a2 \quad\cdots\text{④}
\end{align*}
$\triangle BPQ,\ \triangle ABC$に余弦定理を用いて
\begin{align*}
PQ^2=x^2+y^2-2xy\cos\beta,\qquad 1=1+a^2-2a\cos\beta
\end{align*}
後者から$\cos\beta=\dfrac a2$，④とあわせて
\begin{align*}
PQ^2=x^2+y^2-2xy\cdot\frac a2=x^2+y^2-\frac{a^2}2 \quad\cdots\text{⑤}
\end{align*}
④，$0<x\le1,\ 0<y\le a$のもとで$x^2+y^2$を最小化する．相加相乗平均から$x^2+y^2\ge2xy=a$，等号は$x=y=\sqrt{a/2}$の時．この点が制約をみたすには$\sqrt{a/2}\le1$（$\iff a\le2$，常に成立）かつ$\sqrt{a/2}\le a$（$\iff a\ge1/2$）が必要．

\begin{itemize}
\item $\dfrac12\le a<2$の時：$x=y=\sqrt{a/2}$が制約内にあり，$\min(x^2+y^2)=a$．
\item $0<a<\dfrac12$の時：$\sqrt{a/2}>a$となり$y=\sqrt{a/2}$は$y\le a$に反するので，境界$y=a$（このもとで④から$x=1/2$）で最小となり，$\min(x^2+y^2)=\dfrac14+a^2$．
\end{itemize}
⑤とあわせて
\begin{align*}
\min PQ^2=
\begin{cases}
\dfrac12a^2+\dfrac14 & \left(0<a\le\dfrac12\right)\\[2mm]
-\dfrac12a^2+a & \left(\dfrac12\le a<2\right)
\end{cases}
\end{align*}

\textbf{比較} $1^\circ$の$a^2/2$と$2^\circ$を比較する．$0<a\le1/2$では$2^\circ=a^2/2+1/4>a^2/2=1^\circ$．$1/2\le a<2$では
\begin{align*}
2^\circ-1^\circ=\left(-\frac{a^2}2+a\right)-\frac{a^2}2=a(1-a)
\end{align*}
であり，$1/2\le a<1$で$\ge0$（$1^\circ$が小さい），$1<a<2$で$<0$（$2^\circ$が小さい），$a=1$で等しい．以上より，どちらの場合も$0<a\le1$では$1^\circ$（$PQ^2=a^2/2$）が，$1\le a<2$では$2^\circ$（$PQ^2=-a^2/2+a$）が全体の最小値を与える．したがって，もとめる線分の最小値は
\begin{align*}
\min PQ=
\begin{cases}
\dfrac{\sqrt2}2\,a & (0<a\le1)\\[2mm]
\sqrt{-\dfrac12a^2+a} & (1\le a<2)
\end{cases}
\end{align*}


\newpage

\end{document}
